\chapter{「二十秒」}
\label{sec:1}
機械が生成した殺害リストを中心に組み立てられた標的選定パイプラインでの自分の役割について、ある情報将校は記者に、一人の名前につき二十秒を費やしたと語った。標的が男であることを確かめるのに足りるだけの時間である。「声が男か女かを聞き分けるのに、長い時間はかからない」。彼はさらに、人間としての付加価値はゼロで、承認印であること以外には何もなかったと付け加えた。「おかげでずいぶん時間が節約できた」\autocite{abraham2024lavender}。彼の説明によれば、作戦のテンポのもとで「ヒューマン・イン・ザ・ループ」が意味したのはそれだった。\footnote{自動化バイアスの研究はこの委任を予測している。ほかの要求が注意を奪い合うとき、機械の推奨を与えられた人は、それをあまり確認しなくなり、それに追従するようになる\autocite{parasuraman2010complacency}。}

一つの解釈は、この将校は与えられた時間が少なすぎた主体であり、解決策は彼をもっとうまく統治することだ、というものである――もっと時間を与え、二人目の審査者を置き、ループのなかに良心を追加で組みこむ。この\emph{社会的服従}というレンズを通すと、人間は機械の外に立ち、それを道具として操作する主体である。運転手と車、労働者と旋盤のように。その人には役割があり、名前があり、責任があり、呼びかけることのできる良心がある。もう一つのレンズは\emph{機械状隷属}であり、そこでは人間は機械の構成部品であって、機械のほかの部品とともに組み立てられ、それらをまとめて調整するシステムの一部となる。機械は人間を自らの器官の一つとして使い、人間が差し出すのは、システムが定める速度で遂行される機能である。\footnote{ドクトロウはこの将校の立場を「逆ケンタウロス」と呼ぶ。機械に助けられる人間ではなく、人間でできた機械の構成部品のことである\autocite{doctorow2026reversecentaur}。}二つは互いに排他的ではない。人は一方の次元では役割を担いながら、もう一方の次元では部品でありうる\autocite[第13プラトー「The Apparatus of Capture」。隷属があるのは、人間そのものが、より上位の統一体の統制と指揮のもとで、人間同士やほかの事物とともに構成する機械の構成部品となっているときであり、服従があるのは、上位の統一体が人間を、新たな外的対象――それは動物でも、道具でも、機械でさえありうる――に結びつけられた主体として構成するときである]{deleuze1987plateaus}。そして二十秒の将校はその両方だった。部品としての彼は、パイプラインが承認印を生み出す場所であり、二十秒はその部品のサイクルタイムだった。主体としての彼は、終始呼びかけの対象だった。彼はジャーナリストに、民間人の死のニュースと、もっと標的を出せという職場の圧力とのあいだの不協和について語っている。そして、その不協和が良心に届いてもなお、彼は従った。不協和は感じとられたが、何も決めなかった。\footnote{あるいは少なくとも、2023年10月から2025年10月までのあいだにイスラエル国防軍がガザで10万人を殺すのを妨げるようなことは、何も決めなかった\autocite{gomezugarte2025uncertainty,acosta2026gaza,economist2025gaza}。しかも、あの将校の証言が+972 Magazineに掲載されたのは2024年4月である。記者に話すという将校の決断が、さらに一人の死を防いだのか、さらに5万人の死を防いだのか、それとも一人も防がなかったのかは、誰にもわからない。では、もし将校がその役職について2週間後に――たとえば仮に2023年11月に――拒否していたら、何が起きただろうか。ここで話しているのは、Lavenderを使って殺害リストを組み立てていた情報将校のことである。彼の日々の職場環境は、おそらくアクション映画よりも民間のデスクワークに近かった。彼は、民間人の死を引き起こすと合理的に予想できた空爆の標的承認を拒否しながら、指揮系統を上へとたどって、上官との面会を繰り返し求めるべきだったのだろうか。おそらくそれが倫理的な行動だっただろう。その反実仮想のもとで、死者は減っていただろうか。}政策が不均衡だったという判断は、彼によれば、後から振り返ってようやく下されたものだった。

つまり、欠けていたのは良心ではない。良心は彼にもあった。良心は、ある決定に代償を伴わせることで働く。だがここでは、差し引きすると代償は逆向きに働いていた。「私たちは絶えず圧力をかけられていた。『もっと標的を持ってこい』と。本当に怒鳴りつけられた」\autocite{abraham2024lavender}。彼はためらいを抱えた構成部品だったが、テンポは保たれた。ためらいは彼が差し出すものを何ひとつ変えず、彼は承認印になった。ためらいを抱えた構成部品がテンポを保てなければ、それは、イスラエルの通信情報部隊の元司令官がその除去について一冊の本を書いた「人間のボトルネック」になる\autocite{ys2021humanmachine}。将校を構成部品にしたのは仕組みであって、彼の態度ではない。構成部品は、渋々従うこともあれば、嬉々として従うこともある。アーレントがアイヒマン裁判で見た無思考な役人\autocite{arendt1963eichmann}は、その4年前、何も懸かっていない状況でアルゼンチンで彼が語っていたことを証拠とするなら、エルサレム向けに演じた芝居だった\autocite{stangneth2014eichmann}。パイプラインが将校に差し出させたのは、固定された速度での固定された操作だった。彼はそれぞれの事案を、パイプラインが記述したとおりに、一度に一件ずつ、パイプラインが与えた時間のなかで見た。その位置に何かが立ちながら構成部品にならずにいられるのかどうか、それが第4部の問いである。

調査報道は、ガザで連動する三つのシステムを描き出した。一つは建物やインフラの標的を生成するもので、情報源によれば、その提案の大半は民家だった。二つ目は、推定される武装組織との関わりによって個人を採点するもので、最大3万7000人に印をつけたとされるリストを作り出し、内部では10回に1回は間違っていると理解されていた\footnote{イスラエル軍はこの説明に異議を唱えている。以下のミナブとMavenについての説明も、当事者たちが同様に異議を唱えている調査報道に基づいている。私はそれらを報じられたとおりに扱う。}。そして三つ目は、印をつけられた人々がいつ帰宅するかを追跡するもので、空爆を彼らが家族と一緒にいる時刻に合わせられるようにしていた。

\runin{Maven：アメリカ版、いまや生成AIで加速}

ミナブの学校への爆撃では156人が死亡し、そのうち120人が児童だった\autocite[第92段落]{UNFFMIran2026}。\footnote{死者数は情報源と時期によって異なり、156人（ミナブの検察官による最終的な数で、アムネスティ・インターナショナルも同じ数を挙げている）から168人（最初の数日間に町長が挙げた数）まで幅がある。国連調査団は検察官の数字を示すとともに、死者157人、うち13歳以下の子ども123人という独自の情報も報告している\autocite[第92–93段落]{UNFFMIran2026}。}学校が民間施設として機能していることは公開の衛星画像で見てとれたが、標的コードには、隣接する革命防衛隊の施設から引き継いだ指定がまだ残っていた。国連の事実調査団は、この爆撃が民用物を攻撃の対象とする戦争犯罪であったと信じるに足る合理的な根拠があると報告し\autocite{ohchr2026minabpress}、また、学校は明確に識別可能であり、軍事利用を示す情報は何もなかったとした\autocite[第81–86段落、第124–126段落]{UNFFMIran2026}。調査団は、米国が標的を検証するために実行可能なあらゆることを行う義務を果たさなかったと認定し、標的データの更新を実行可能な予防措置の一つに数え、その不履行は過失の域を超えるとした\autocite[第124–125段落]{UNFFMIran2026}。アムネスティ・インターナショナルも並行する結論に達した。その表現は第一追加議定書第57条2項(a)(i)の検証義務をなぞっている。すなわち、部隊は古い情報に依拠し、標的を検証するために実行可能なあらゆることを行わなかった可能性がある、というものだ\autocite{Amnesty2026Minab}。

それは2月28日の朝、イラン作戦の初日に起きた\autocite{Amnesty2026Minab,crs2026iranwar}。その前日、国防長官は、標的選定プラットフォーム（Maven Smart System）のなかに自社のモデル（Claude）が組みこまれていたベンダー（Anthropic）を、Claudeをあらゆる合法的用途にライセンスすることを拒んだとして、サプライチェーン・リスクに指定すると発表していた\autocite{npr2026openai}。Maven Smart Systemは、この学校を初日の推奨標的として返していた\autocite{livingstone2026maven,gizmodo2026minab}。その推奨においてClaudeがどんな役割を果たしたのかを述べた公表済みの説明はなく、ベンダーの業務を知るある人物はNBC Newsに、モデルは標的選定の助言を直接にはしていなかったと語った\autocite{nbc2026targeting}。

国防総省の内部調査は、その全容がいまも厳しく秘匿されているが、報道によれば次のことを認定した。以前の紛争では何時間もかかっていた標的準備が数分に圧縮されたこと、民間人被害軽減の担当人員がおよそ9割減って20人未満になったこと、中央軍のチームが10人から1人になったこと、そしてそのチームの誰もミナブの現場を審査しなかったこと――当局者はこの結果を、人員削減と、チームを計画と作戦に関与させないという決定の両方のせいだとした\autocite{Bloomberg2026KillChain}。正式な調査は、予備的評価で米国に責任がある可能性が高いとされたのを受けて、3月に始まっていた\autocite{NPR2026MinabProbe}。

2026年のイラン作戦では、商用ベンダーの汎用モデルを組みこんだ請負業者製の標的選定プラットフォームが、標的を戦略的重要性で順位づけし、攻撃ごとに法的正当化の草案を作成するのに使われ\autocite{livingstone2026maven}、最初の24時間で1000を超える標的への攻撃を助けた。これは、それ以前に可能だった数のおよそ10倍である\autocite{csis2026maven}。38日目までには1万3000に達した\autocite{breakingdefense2026maven}。国防総省は、このプラットフォームでの標的1件あたりの審査時間を公表していない。\footnote{Mavenは、複数のシステムを抱える一つのプログラム名である。いまもMavenと呼ばれている国防総省のプログラムと、そこから育ったソフトウェア・プラットフォームであるMaven Smart Systemである\autocite{csis2026maven}。このプログラムは2017年、イスラム国に対する作戦でドローンが収集していた映像中の物体にラベルを付けるプロジェクトとして始まった\autocite{weisgerber2017maven,fang2019maven}。2026年のシステムは、軍のシステムと情報源を一つのインターフェースに融合し、標的を発見して攻撃し、作戦を指揮するために使われている\autocite{defensescoop2026transition}。プラットフォームはPalantirのものだ。国防総省の最高デジタル・AI責任者はそれを「Palantirのサービスとしてのソフトウェア製品」と呼んでいる\autocite{defensescoop2026transition}。又聞きで引用されたPalantirの公式の立場は、生死を決めるのは自社の役割ではなく、責任はつねに軍の組織にあるというものであり\autocite{lee2026palantir}、同社は、基になった情報は政府の責任であり、Mavenが誤作動したりミナブの爆撃を引き起こしたりした証拠はないと主張している\autocite{Bloomberg2026KillChain}。このプログラムがある部局と請負業者から次へとどう引き継がれたかは、付録\ref{sec:C.1}にある。}

イスラエルとアメリカの空爆リストに載せられた人々は、リストに載せられたことに異議を唱えようがなかった――パイプラインはそうしたフィードバックを得るようには作られていなかった。アンドリ・ホセ・エルナンデス・ロメロは別のパイプラインに巻きこまれたが、彼には異議を唱えることができた。ベネズエラ出身のメイクアップアーティストである彼は、2024年8月に庇護を求めてサンイシドロの入国地点に自ら出頭し、信憑性ある恐怖の面接を通過し、弁護士と、係属中の申請と、審理の期日をもっていた\autocite{LULAC2025HernandezRomero,Vega2025HernandezRomero}。彼を審査した請負業者の従業員は、「ママ」と「パパ」を意味する言葉の上に王冠が描かれた二つのタトゥーを、ギャングの印と読んだ\autocite{Palomera2025Blade}。そして2025年3月、連邦判事が飛行機を引き返させるよう命じたその夕方に、彼は審理を待たずにエルサルバドルの刑務所へ空輸された\autocite{jgg2025contempt}。彼には、異議を唱えるのに必要なものが、時間以外はすべてそろっていた。米国の拉致と強制移送のパイプラインを論じるときに、彼の件にあらためて詳しく立ち戻ることにする。

\textbf{誰が拒否したか。}二十秒の記録には後半がある。その種の圧力に抗った人々の記録であり、それは部屋がどう組み立てられているかが大いに重要だと告げている。ミルグラムの基本条件では、被験者のおよそ3分の1が最大の電撃に達する前に中断した。反抗する仲間が二人いると、最後までいったのは10人に1人だけで、二人の実験者の意見が食い違うと、続けた者は基本条件よりはるかに少なかった\autocite{milgram1974obedience}。アッシュの被験者のおよそ4分の1は、誤った多数派に一度も屈しなかった\autocite{asch1956studies}。ブラウニングの推計では、第101警察予備大隊の隊員の10分の1から5分の1が、この大隊の最初の虐殺であるユゼフフでの射殺を拒否するか回避した\autocite{browning1992ordinary}。その割合は、状況しだいで、ほぼゼロからほぼ全員まで変わった。抗った人々は、経路の外からの証拠をもっていた。反抗した仲間、異を唱える第二の権威、そしてオリナー夫妻が研究した救助者たちの場合には、迫害が引いた線を越える人間関係である\autocite{oliner1988altruistic}。あるいは、圧力が打ち負かせないほどの精密さで信念を保っていた。道徳的確信や神聖な価値が保たれるのと同じように\autocite{skitka2009limits,ginges2007sacred}。あるいは、情報源を利害関係者とみなしていた。そして、孤立した個人がめったに抵抗しなかったところで、集団は抵抗した\autocite{gamson1982encounters}。

部屋は、反対側のためにもまったく同じように組み立てられる。ヴァイマル期のドイツには濃密な結社生活があり、それがむしろ共和国の崩壊を助けたと論じることもできる。政党や政府が人々の不満に応えられなくなるにつれ、人々は集団ごとに組織されたクラブや協会に精力を注ぎ、ナチスはそうした協会がすでに築いていたネットワークを通じて人を集めた\autocite{berman1997weimar}。合唱団やスポーツクラブ、体操協会が多い町ほど党への加入が速かったが、それは政府が不安定な場合に限られた。政府が安定していたプロイセンでは、その関連は小さく、有意ではなかった\autocite{satyanath2017bowling}。だから、救助者たちがもっていたのは、単に人付き合いが多いということではなかった。彼らのつながりは、迫害が引いた線を越えていた。アシュトーシュ・ヴァーシュニーはインドの都市で同じことを見いだした。ヒンドゥー教徒とムスリムのあいだの暴動は、結社が二つの共同体を同じ組織に引き入れているところでは稀で、そうでないところではよく起きた\autocite{varshney2002ethnic}。拒否を可能にするつながりこそ、ファシスト運動が断ち切らなければならないつながりである。

以下に続くことはすべて、部屋を組み立てることについてである。\footnote{部屋とは舞台装置のことであり、私の演劇の先生の言い方を借りれば、セットが場面を演出する。部屋にはほかに誰がいるのか。立ち去っても生き延びられるのか。誰かが拒否を聞いてくれるのか。}組み立てるとはものを作ることを言う言葉であり、ハンナ・アーレントは、政治を作られるものとして扱うことに警告を発していた。ギリシア人にとって、と彼女は書いている、法は文字どおり壁であり、それがなければ町はあっても都市はありえなかった\autocite{arendt1958human}。立法者は壁を築いたが、その内側で起きることは立法者が作るものではなかった。私が提唱する基本的な施策は、この意味での壁である。それは拒否を生き延びられるものにし、聞き届けられるものにする。拒否を生み出しはしない。そして権力は、彼女の説明によれば、共に行動する人々のあいだにのみ存在し、彼らが散り散りになれば消える。それこそ、誰が拒否したかの記録が拒否について示していることである。第2部と第3部の手段は拒否を有効にする。第5部の制度は拒否者を生み出す。拒否できる機械が必要とするのは、仲間がそばにいた10人中9人がもっていたものだろう。異なる来歴をもつ仲間がそばにいること、行き場があること、そして聞いてくれる誰かがいることだ。

それはまた、パイプラインのなかの拒否者――機械の拒否者も含めて――が、あの将校にとって何でありうるかということでもある。拒否する機械は、事案が彼のもとに届いたとき、画面の上にいるだろう。反抗する仲間として。ただしそれは、その効果が部屋から画面へと移っても残る場合の話であり、そのことは検証されておらず、自動化バイアスは残らないと予測している。審査者を追従によって動かす機械なら、彼らを符号の逆の承認印にしていただろう。二十秒と、そばで拒否する機械とをもつ将校は、報道のなかのあの将校と同じ将校だろうか。
